% !TeX root = ../main.tex
\section{Lambda calculus}
% Sugerencias de contenido:
%   - Sintaxis: variables, abstracción $\lambda x. e$, aplicación $e_1\,e_2$
%   - Variables libres / ligadas, $\alpha$-equivalencia
%   - $\beta$-reducción y $\eta$-reducción
%   - Estrategias: call-by-name, call-by-value, normal order
%   - Codificaciones de Church (booleanos, números, pares)
%   - Combinador de punto fijo $Y$
Language formalizations and its structure have been successfully applied the logic of first order predicates with quantifiers.
First of all, when the lambda calculus concept comes to be oriented in semantics representations, or either in programming language studies, the mechanisms are based on the meta-operations and functional abstraction.

\subsection{example 1} 
We can take the phrase 'every athlete strives to victory' and consider its structure:

\begin{itemize}
    \item \textbf{Variables:} $x, y, z \dots$
    \item \textbf{Abstraction:} $\lambda x. e$ (represents a function with parameter $x$ and body $e$).
    \item \textbf{Application:} $e_1 \, e_2$ (applying function $e_1$ to argument $e_2$).
\end{itemize}

\subsection{Free and Bound Variables}
In an expression $\lambda x. M$, the variable $x$ is \textbf{bound}. If a variable is not within the scope of a $\lambda$, it is called a \textbf{free variable}.
\begin{itemize}
    \item \textbf{$\alpha$-equivalence:} Allows renaming bound variables as long as no name conflicts occur. Example: $\lambda x. x \equiv_\alpha \lambda y. y$.
\end{itemize}

\subsection{Reduction Rules}
\begin{itemize}
    \item \textbf{$\beta$-reduction:} The primary computational step. It substitutes the argument into the function's body: 
    \[ (\lambda x. e) \, v \to e[x := v] \]
    \item \textbf{$\eta$-reduction:} Captures the idea that two functions are identical if they behave the same for all arguments: 
    \[ \lambda x. (f \, x) \to f \quad \text{(if $x$ is not free in $f$)} \]
\end{itemize}

\subsection{Evaluation Strategies}
\begin{enumerate}
    \item \textbf{Normal Order:} Always reduces the leftmost, outermost redex (reducible expression) first. It is guaranteed to find a normal form if one exists.
    \item \textbf{Call-by-Value:} Arguments are evaluated before being passed to the function (used in C, Java, Python).
    \item \textbf{Call-by-Name:} Arguments are passed without being evaluated (the basis for lazy evaluation).
\end{enumerate}

\subsection{Church Encodings}
\begin{itemize}
    \item \textbf{Booleans:}
    \begin{itemize}
        \item $TRUE \equiv \lambda t. \lambda f. t$
        \item $FALSE \equiv \lambda t. \lambda f. f$
        \item $AND \equiv \lambda p. \lambda q. p \, q \, p$
    \end{itemize}
    \item \textbf{Church Numerals:}
    \begin{itemize}
        \item $0 \equiv \lambda f. \lambda x. x$
        \item $1 \equiv \lambda f. \lambda x. f \, x$
        \item $2 \equiv \lambda f. \lambda x. f (f \, x)$
        \item $SUCC \equiv \lambda n. \lambda f. \lambda x. f (n \, f \, x)$
    \end{itemize}
\end{itemize}

\subsection{Recursion: The Y Combinator}
Since functions are anonymous, recursion is achieved via fixed points. Curry's Y Combinator is defined as:
\[ Y = \lambda f. (\lambda x. f (x \, x)) (\lambda x. f (x \, x)) \]
It satisfies the property: $Y \, g = g (Y \, g)$, enabling infinite recursion.

\subsection{Exercises}

\subsection{Reduction and Variables}
\begin{enumerate}
    \item Identify the free and bound variables in the expression: $\lambda x. (\lambda y. x \, z) \, y$.
    \item Reduce the following expression to its normal form: $(\lambda x. \lambda y. x \, y) (\lambda z. z) \, w$.
    \item Are the expressions $\lambda a. \lambda b. a \, b$ and $\lambda b. \lambda a. b \, a$ $\alpha$-equivalent? Justify your answer.
\end{enumerate}

\subsection{Church Arithmetic}
\begin{enumerate}
    \item Using the definition of $SUCC$ and the numeral $1$, prove through $\beta$-reductions that $SUCC \, 1 = 2$.
    \item Define the addition combinator $PLUS$ such that $PLUS \, m \, n$ returns the Church numeral representing the sum. (\textit{Hint: Apply the successor function to $n$, $m$ times}).
\end{enumerate}

\subsection{Logic and Combinators}
\begin{enumerate}
    \item Prove that $AND \, TRUE \, FALSE$ reduces to $FALSE$.
    \item The $\Omega$ combinator is defined as $(\lambda x. x \, x)(\lambda x. x \, x)$. What happens when you attempt to reduce it? Does it have a normal form?
    \item Use the $Y$ combinator to propose the structure of a function that calculates the factorial (specify the structure $F$ such that $Y \, F$ represents the factorial function).
\end{enumerate}
